\documentclass[../../../include/open-logic-section]{subfiles}

\begin{document}

\olfileid{sth}{choice}{vitali}
\olsection{ضمیمہ: ویتالی دا پیراڈوکس}

باناخ–تارسکی دی بناوٹ قبول کرنی ٹھیک اے یا نہیں، ایہ چنگی طرح سمجھن لئی سانوں اوہدا \emph{ثبوت} ویکھنا چاہیدا اے۔ افسوس، اوہدے لئی اینی الجبرا چاہیدی اے جِنّی اسیں ایتھے نہیں دس سکدے۔ پر ہِیٹھلے نتیجے اُتے دھیان دے کے اسیں کجھ مختصر گلاں دس سکدے آں جیہڑیاں باناخ–تارسکی دے ثبوت دی کجھ سمجھ دے سکدیاں نیں۔\footnote{کِتے ودھ تفصیل لئی \cite{Weston2003} یا \cite{Wagon2016} ویکھو۔}

\begin{thm}[ویتالی دا پیراڈوکس ($\ZFC$ وچ)]\ollabel{vitaliparadox}
کسے وی دائرے نوں گِنن جوگ گنتی دے ٹکڑیاں وچ ونڈیا جا سکدا اے۔ اوہ ٹکڑے فیر (گھما کے تے تھاں بدل کے) جوڑے جا سکدے نیں تاں جو اوس دائرے دیاں دو نقلاں بن جان۔
\end{thm}

ویتالی دا پیراڈوکس باناخ–تارسکی دے پیراڈوکس نالوں ثابت کرنا کِتے سوکھا اے۔ اسیں اینوں «ویتالی دا پیراڈوکس» آکھیا اے، کیوں جے ایہ ویتالی دی \citeyear{Vitali1905} وچ اک ناپن جوگ نہ ہون والے سیٹ دی بناوٹ توں نکلدا اے۔ پر ویتالی تے باناخ–تارسکی دے پیراڈوکس دے ثبوتاں وچ سیٹ نظریے والے پاسے بڑے ملدے نیں۔ نتیجیاں وچ اصلی فرق بس ایہ اے کہ باناخ–تارسکی \emph{محدود} گنتی دے ٹکڑیاں والی ونڈ ویکھدا اے، جدکہ ویتالی دا پیراڈوکس \emph{گِنن جوگ لامتناہی} گنتی دے ٹکڑیاں والی ونڈ ویکھدا اے۔ جیویں \citet{Weston2003} آکھدا اے، ویتالی دا پیراڈوکس «باناخ–تارسکی دے پیراڈوکس جِنّا حیران کن ہرگز نہیں، پر ایہ وکھاندا اے کہ ہندسی پیراڈوکس ‘سادیاں’ صورتاں وچ وی ہو سکدے نیں»۔

ویتالی دا پیراڈوکس اک دو جہتی شکل، دائرے، بارے اے۔ سو اسیں سطح $\Real^2$ اُتے کم کراں گے۔ $\rotationsgroup$ نقطیاں دے مبدأ دوالے (گھڑی دے کانٹیاں دی سمت) گھماؤواں دا سیٹ ہووے، جیہناں دے زاویے \emph{ناطق} تے $[0,2\pi)$ وچ ریڈیئن قدراں نیں۔ $\rotationsgroup$ بارے الجبرا دیاں کجھ گلاں ایہ نیں (جے نتیجے دا بیان سمجھ نہ آوے تاں ثبوت اوہدا مطلب صاف کر دے گا):

\begin{lem}\ollabel{rotationsgroupabelian}
$\rotationsgroup$ فنکشناں دی ترکیب تحت اک ابیلی {گروپ} بندا اے۔
\end{lem}

\begin{proof}
$0_{\rotationsgroup}$ نال $0$ ریڈیئن دا گھماؤ لکھدے ہوئیاں، ایہ $\rotationsgroup$ دا شناختی رکن اے، کیوں جے
$\comp{0_{\rotationsgroup}}{\rho} = \comp{\rho}{0_{\rotationsgroup}} =
\rho$ کسے وی $\rho \in \rotationsgroup$ لئی۔

ہر رکن دا اک اُلٹ اے۔ جے $\rho \in \rotationsgroup$، $r$ ریڈیئن گھماؤندا اے، تاں $\rho^{-1} \in \rotationsgroup$، $2\pi - r$ ریڈیئن گھماؤندا اے، سو $\rho \circ \rho^{-1} = 0_\rotationsgroup$۔

ترکیب وچ شرکتیت اے۔ $\comp{\rho}{(\comp{\sigma}{\tau})} =
\comp{(\comp{\rho}{\sigma})}{\tau}$ کسے وی $\rho, \sigma, \tau \in
\rotationsgroup$ لئی۔

ترکیب تبادلی اے۔ $\comp{\rho}{\sigma} =
\comp{\sigma}{\rho}$ کسے وی $\rho, \sigma \in \rotationsgroup$ لئی۔
\end{proof}

اصل وچ اسیں اپنے گروپ $\rotationsgroup$ نوں دو ادھّاں وچ ونڈ سکدے آں، تے فیر کسے وی اک ادھ نال سارا گروپ فیر حاصل کر سکدے آں۔

\begin{lem}\ollabel{disjointgroup}
$\rotationsgroup$ دی دو وکھرے سیٹاں، $\rotationsgroup_{1}$ تے $\rotationsgroup_{2}$، وچ ونڈ اے، تے دوویں $\rotationsgroup$ دی بنیاد نیں۔
\end{lem}

\begin{proof}
$\rotationsgroup_{1}$ وچ اوہ گھماؤ ہون جیہناں دیاں ناطق ریڈیئن قدراں $[0, \pi)$ وچ نیں۔ $\rotationsgroup_{2} = \rotationsgroup
\setminus \rotationsgroup_1$ ہووے۔ بنیادی الجبرا توں
$\Setabs{\comp{\rho}{\rho}}{\rho \in \rotationsgroup_1} =
\rotationsgroup$۔ $\rotationsgroup_2$ لئی وی ایہو جیہا نتیجہ مل سکدا اے۔
\end{proof}

اسیں گروپاں بارے ایس گل نوں \olref{vitaliparadox} ثابت کرن لئی ورتاں گے۔ $\onesphere$ اکائی دائرہ ہووے، یعنی اوہناں نقطیاں دا سیٹ جیہڑے سطح دے مبدأ توں این ٹھیک~$1$ اکائی دور نیں، یعنی
$\Setabs{\tuple{r,s} \in \Real^2}{\sqrt{r^2+s^2}=1}$۔ اسیں $\onesphere$ نوں حصیاں وچ ونڈاں گے، $\onesphere$ اُتے ہِیٹھلا تعلق ویکھ کے:
\[
	r \sim s \emph{ اُدوں تے صرف اُدوں جدوں }(\exists \rho \in \rotationsgroup)\rho(r) = s.
\]
یعنی $\onesphere$ دے دو نقطے ایس تعلق نال جڑے نیں اُدوں تے صرف اُدوں جدوں مبدأ دوالے ناطق قدر والے گھماؤ نال اک توں دوجے تک پہنچ سکدے او۔ ہُن ایہ گل حیرانی والی نہیں:

\begin{lem}
$\sim$ اک ہم ارزی دا تعلق اے۔
\end{lem}

\begin{proof}
\olref{rotationsgroupabelian} ورت کے ایہ گل صاف اے۔
\end{proof}

ہُن اسیں چُون ورت کے اک سیٹ، $C$، حاصل کردے آں جیہدے وچ $\onesphere$ دی $\sim$ تحت ہر ہم ارزی جماعت وچوں این ٹھیک اک رکن اے۔ یعنی اسیں ہم ارزی جماعتاں دے سیٹ اُتے اک چُون والا فنکشن $f$ لیندے آں۔\footnote{کیوں جے $\rotationsgroup$ !!{enumerable} اے، $E$ دا ہر !!{element} !!{enumerable} اے۔ کیوں جے $\onesphere$ !!{nonenumerable} اے، \olref{vitalicover} تے \olref[card-arithmetic][simp]{kappaunionkappasize} توں نکلدا اے کہ $E$ !!{nonenumerable} اے۔ سو ایتھے \emph{گِنن جوگ نہ ہون والے} سیٹ لئی چُون ورتی اے۔}
\[
	E = \Setabs{\equivrep{r}{\sim}}{r \in \onesphere},
\]
تے $C = \ran{f}$ ہووے۔ ہر گھماؤ $\rho \in \rotationsgroup$ لئی، سیٹ $\funimage{\rho}{C}$ اوہناں نقطیاں توں بندا اے جیہڑے گھماؤ $\rho$ نوں $C$ دے ہر نقطے اُتے لاؤن نال ملدے نیں۔ اگلے دو نتیجے وکھاندے نیں کہ ایہ سیٹ دائرے نوں پورا ڈھکدے نیں تے اوہناں وچ کوئی سانجھا نقطہ نہیں۔

\begin{lem}\ollabel{vitalicover}
$\onesphere = \bigcup_{\rho \in \rotationsgroup} \funimage{\rho}{C}$۔
\end{lem}

\begin{proof}
$s \in \onesphere$ مقرر کرو۔ کوئی $r \in C$ اے جیہدے لئی $r \in
\equivrep{s}{\sim}$، یعنی $r \sim s$، یعنی $\rho(r) = s$ کسے $\rho \in \rotationsgroup$ لئی۔
\end{proof}

\begin{lem}\ollabel{vitalinooverlap}
جے $\rho_1 \neq \rho_2$، تاں $\funimage{\rho_1}{C} \cap \funimage{\rho_2}{C} = \emptyset$۔
\end{lem}

\begin{proof}
منّ لو $s \in \funimage{\rho_{1}}{C} \cap \funimage{\rho_{2}}{C}$۔ سو
$s = \rho_{1}(r_{1}) = \rho_{2}(r_{2})$ کجھ $r_{1}, r_{2} \in C$ لئی۔
ایس لئی $\rho^{-1}_2(\rho_1(r_1)) = r_2$، تے
$\comp{\rho_1}{\rho^{-1}_2} \in \rotationsgroup$، سو $r_{1} \sim
r_{2}$۔ سو $r_1 = r_2$، کیوں جے $C$، $\sim$ تحت ہر ہم ارزی جماعت وچوں این ٹھیک اک رکن چُندا اے۔ سو $s = \rho_1(r_1) = \rho_2(r_1)$، تے ایس لئی $\rho_1 = \rho_2$۔
\end{proof}

ہُن اسیں الجبرا دیاں پہلاں والیاں گلاں اپنے دائرے اُتے لاؤندے آں۔

\begin{lem}\ollabel{pseudobanachtarski}
$\onesphere$ دی دو وکھرے سیٹاں، $D_{1}$ تے $D_{2}$، وچ ونڈ اے، ایویں کہ $D_{1}$ نوں گِنن جوگ گنتی دے سیٹاں وچ ونڈیا جا سکدا اے جیہناں نوں گھما کے $\onesphere$ دی اک نقل بنائی جا سکدی اے (تے $D_{2}$ لئی وی ایہو گل اے)۔
\end{lem}

\begin{proof}
\olref{disjointgroup} والے $\rotationsgroup_{1}$ تے $\rotationsgroup_{2}$ ورت کے ایویں لئیے:
\begin{align*}
	D_{1} &= \bigcup_{\rho \in \rotationsgroup_1} \funimage{\rho}{C} &
	D_{2} &= \bigcup_{\rho \in \rotationsgroup_2} \funimage{\rho}{C}
\end{align*}
\olref{vitalicover} مطابق ایہ $\onesphere$ دی ونڈ اے، تے $D_1$ تے $D_2$، \olref{vitalinooverlap} مطابق وکھرے نیں۔ بناوٹ مطابق $D_1$ نوں گِنن جوگ گنتی دے سیٹاں وچ ونڈیا جا سکدا اے، یعنی $\funimage{\rho}{C}$ ہر $\rho \in R_1$ لئی۔ تے ایہناں نوں گھما کے $\onesphere$ دی نقل بنائی جا سکدی اے، کیوں جے $\onesphere = \bigcup_{\rho \in
\rotationsgroup}\funimage{\rho}{C} = \bigcup_{\rho \in
\rotationsgroup_1}\funimage{(\comp{\rho}{\rho})}{C}$، \olref{disjointgroup} تے \olref{vitalicover} مطابق۔ اوہی دلیل $D_2$ لئی چل دی اے۔ \end{proof}\noindent ایس توں ویتالی دا پیراڈوکس جھٹ نکلدا اے۔ کیوں جے اسیں $\onesphere$ دیاں \emph{دو} نقلاں $\onesphere$ توں بنا سکدے آں۔ بس اینوں گِنن جوگ گنتی دے ٹکڑیاں (وکھ وکھ $\funimage{\rho}{C}$ والے سیٹاں) وچ ونڈو تے فیر اوہناں دی شکل قائم رکھدیاں ہلاؤ ($D_1$ دا ہر ٹکڑا بس گھماؤ، تے $D_2$ دے ہر ٹکڑے دی پہلاں تھاں بدلو تے فیر گھماؤ)۔

ثبوت دا طریقہ فیر چیتے کریے۔ اسیں سطح اُتے گھماؤواں دے گروپ بارے الجبرا دیاں کجھ گلاں توں شروع کیتا۔ ایس گروپ نال اسیں $\onesphere$ نوں ہم ارزی جماعتاں وچ ونڈیا۔ فیر چُون ورت کے ہر جماعت وچوں رکن چُنے تے ایس طرح «پیراڈوکس» تک پہنچے۔

باناخ–تارسکی ثابت کرن لئی این اوہی طریقہ ورتدے آں۔ وڈا فرق ایہ اے کہ باناخ–تارسکی دے ثبوت وچ ورتیاں الجبرا دیاں گلاں ویتالی دے پیراڈوکس دے ثبوت والیاں نالوں کِتے گنجل دار نیں۔ پر اوہناں الجبرا دیاں گلاں دا چُون نال کوئی تعلق نہیں۔ اسیں اوہناں دا مختصر خلاصہ دیواں گے۔

باناخ–تارسکی ثابت کرن لئی اسیں \olref{disjointgroup} ورگی اک گل ثابت کرن توں شروع کردے آں۔ کسے وی \emph{آزاد گروپ} نوں چار ٹکڑیاں وچ ونڈیا جا سکدا اے۔ سہج سمجھ مطابق اسیں ایہناں نوں «اگے پچھے ہلا» کے سارے گروپ دیاں دو نقلاں فیر حاصل کر سکدے آں۔\footnote{ایہ گل کہ اسیں \emph{چار} ٹکڑے ورت سکدے آں، \cite{Robinson1947} دی اے۔ اک حالیہ ثبوت لئی \citet[Theorem 5.2]{Wagon2016} ویکھو۔ گروپ دے ٹکڑیاں نوں «ہلاؤن» والی گل آکھدیاں اسیں \citet[p.~3]{Weston2003} دے پچھے چلدے آں۔} فیر اسیں وکھاندے آں کہ $\Real^3$ دے مبدأ دوالے دو خاص گھماؤ ورت کے گھماؤواں دا اک آزاد گروپ، $F$، پیدا کر سکدے آں۔\footnote{\citet[Theorem 2.1]{Wagon2016} ویکھو۔} (ہُن تک چُون نہیں ورتی۔) ہُن اسیں گولے دی سطح دے نقطیاں نوں «ملدے» سمجھدے آں اُدوں تے صرف اُدوں جدوں اک توں دوجا~$F$ دے کسے گھماؤ نال مل سکدا اے۔ فیر اسیں «ملدے» نقطیاں دی ہر ہم ارزی جماعت وچوں این ٹھیک اک نقطہ چُنن لئی \emph{چُون ورتدے آں}۔ $F$ دی اپنی ونڈ نوں گولے دی سطح اُتے لا کے، جیویں \olref{pseudobanachtarski} وچ، اسیں سطح نوں چار ٹکڑیاں وچ ونڈدے آں جیہناں نوں «اگے پچھے ہلا» کے گولے دی سطح دیاں دو نقلاں بنا سکدے آں۔ تے ایس توں ایہ ثابت ہُندا اے \citep{Hausdorff1914}:

\begin{thm}[ہاؤسڈورف دا پیراڈوکس ($\ZFC$ وچ)]
کسے وی گولے دی سطح نوں محدود گنتی دے ٹکڑیاں وچ ونڈیا جا سکدا اے۔ اوہ ٹکڑے فیر (گھما کے تے تھاں بدل کے) جوڑے جا سکدے نیں تاں جو اوس گولے دیاں دو وکھریاں نقلاں بن جان۔
\end{thm}

پورا باناخ–تارسکی قضیہ حاصل کرن لئی (جیہڑا گولے دی سطح دے نال اوہدے اندرلے حصے بارے وی اے) الجبرا دیاں دو کُ ہور ترکیباں چاہیدیاں نیں۔ پر سِدھی گل ایہ اے کہ الجبرا دا اصل کم ہو چکیا اے۔ ایس لئی ویسٹن لکھدا اے:
\begin{quote}
  [\ldots] آزاد گروپاں والا نتیجہ باناخ–تارسکی دے پیراڈوکس دے ثبوت دا \emph{ضروری قدم} اے۔ ایس نظر توں باناخ–تارسکی دا پیراڈوکس $\Real^3$ بارے گھٹ، تے گروپ دی پیچیدگی بارے ودھ اے [یعنی $\Real^3$ وچ تھاں بدلن تے گھماؤن دا گروپ]۔ \cite[p.~16]{Weston2003}
\end{quote}
یعنی اسیں \emph{محدود} گنتی دے ٹکڑیاں والی ونڈ دے سکدے آں (جیویں باناخ–تارسکی وچ)، یا \emph{گِنن جوگ لامتناہی} گنتی دے ٹکڑیاں والی ونڈ (جیویں ویتالی دے پیراڈوکس وچ)، ایہ دو یا تِن جہتاں وچ کم کرن بارے گروپ نظریے دیاں کجھ گلاں اُتے ٹکدا اے۔

منّنا پئے گا کہ ایہ آخری گل \olref[banach]{sec} دے آخر والے لطیفے نوں ذرا خراب کردی اے۔ دو جہتی ہون کرکے «باناخ–تارسکی» نوں گِنن جوگ \emph{لامتناہی} گنتی دے ٹکڑیاں وچ ونڈنا پئے گا، جے اوہ ٹکڑے فیر جوڑ کے «باناخ–تارسکی باناخ–تارسکی» بنانا ہووے۔ لطیفہ ٹھیک کرن لئی تِن جہتاں وچ لکھنا پئے گا۔ اسیں ایہ پڑھن والے لئی مشق وجوں چھڈدے آں۔

اک آخری گل۔ \olref[banach]{sec} وچ اسیں دسیا سی کہ گولے توں ملے «ٹکڑے» \emph{ناپن جوگ} نہیں ہو سکدے، پر نقطیاں دے ایہو جیہے «لامتناہی کھلارے» ہون گے جنہاں دی تصویر ذہن وچ نہیں بن سکدی۔ ایہو گل \olref{pseudobanachtarski} حاصل کرن لئی چُون دی ساڈی ورتوں اُتے وی لاگو ہُندی اے۔ تے ایس سارے دی وضاحت کرنی چنگی اے۔

فیر سانوں کجھ پچھوکڑ دسنا پئے گا (پر ایہ \emph{بس} اک خاکہ اے۔ تسی \emph{ناپ} بارے کسے درسی کتاب دا حصہ ویکھ سکدے او)۔ سیٹ $X$ لئی اک ناپ تعریف کرن دا مطلب اے کہ اک قدر $\mu(E) \in
\Real$ ہر $E$ نوں دِتی جائے، جیہڑا کسے «$\sigma$-الجبرا» وچ اے جیہڑا $X$ اُتے اے۔ ایتھے تفصیل ضروری نہیں، بس فنکشن $\mu$ نوں گِنن جوگ جوڑ دے اصول اُتے پورا اترنا چاہیدا اے۔ وکھرے سیٹاں دی گِنن جوگ گنتی دے اتحاد دا ناپ اوہناں دے وکھ وکھ ناپاں دا جوڑ اے، یعنی
$\mu(\bigcup_{n < \omega} X_n) = \sum_{n < \omega}\mu(X_n)$ جدوں $X_n$ والے سیٹ وکھرے ہون۔ کسے سیٹ نوں «ناپن جوگ نہیں» آکھن دا مطلب اے کہ اوہنوں کوئی ڈھکواں ناپ نہیں دِتا جا سکدا۔ ہُن پہلاں والا $\rotationsgroup$ ورتدے آں:

\begin{cor}[ویتالی]
$\mu$ اک ایہو جیہا ناپ ہووے کہ $\mu(\onesphere) = 1$، تے $\mu(X) = \mu(Y)$ جے $X$ تے $Y$ شکل تے وڈائی وچ برابر نیں۔ تاں
$\funimage{\rho}{C}$ ناپن جوگ نہیں، سارے $\rho \in
\rotationsgroup$ لئی۔
\end{cor}

\begin{proof}
ٹکراؤ کڈھن لئی اُلٹ منّ لو۔ سو $\mu(\funimage{\sigma}{C}) =
r$ ہووے کسے $\sigma \in \rotationsgroup$ تے کسے $r \in \Real$ لئی۔ کسے وی $\rho \in C$ لئی، $\funimage{\rho}{C}$ تے $\funimage{\sigma}{C}$ شکل تے وڈائی وچ برابر نیں، تے ایس لئی $\mu(\funimage{\rho}{C}) = r$ کسے وی $\rho \in
C$ لئی۔ \olref{vitalicover} تے \olref{vitalinooverlap} مطابق، $\onesphere =
\bigcup_{\rho \in \rotationsgroup}\funimage{\rho}{C}$ جوڑی جوڑی کرکے وکھرے سیٹاں دی گِنن جوگ گنتی دا اتحاد اے۔ سو گِنن جوگ جوڑ دا اصول دسدا اے کہ $\mu(\onesphere) = 1$ ہر $\funimage{\rho}{C}$ دے ناپاں دا جوڑ اے، یعنی:
\[
	1 = \mu(\onesphere) = \sum_{\rho \in \rotationsgroup}\mu(\funimage{\rho}{C}) = \sum_{\rho \in \rotationsgroup}r
\]
پر جے $r = 0$، تاں $\sum_{\rho \in \rotationsgroup}r = 0$، تے جے $r
> 0$، تاں $\sum_{\rho \in \rotationsgroup}r = \infty$۔
\end{proof}

\end{document}
